\section{系统架构}
\label{sec:model}

\subsection{系统概述}

系统采用协调者-代理（Coordinator-Agent）架构，由中心化的\texttt{Coordinator}编排四个专业化Agent按顺序执行流水线。数据持久化和Agent间通信通过\texttt{SharedWorkspace}（共享工作区）管理，它为聚类和摘要提供按主题隔离的存储，同时全局共享文献和嵌入向量。

图~\ref{fig:architecture}展示了系统的整体架构。

\begin{figure}[htbp]
    \centering
    \begin{tikzpicture}[
        box/.style={rectangle, draw, rounded corners, minimum width=2.6cm, minimum height=1cm, align=center, font=\small},
        agent/.style={rectangle, draw, rounded corners, fill=blue!10, minimum width=2.2cm, minimum height=1cm, align=center, font=\small},
        arr/.style={-{Stealth[length=2.5mm]}, thick},
        darr/.style={-{Stealth[length=2.5mm]}, thick, dashed}
    ]
        % 用户输入
        \node[box, fill=yellow!20] (input) at (3,5) {用户输入\\研究主题};

        % 协调者
        \node[box, fill=red!10] (coord) at (3,3.5) {Coordinator\\流水线编排};

        % RQ管理器
        \node[box, fill=orange!10] (rq) at (8,3.5) {RQ Manager\\层级研究问题};

        % 四个Agent（均匀排列）
        \node[agent] (search)  at (0,1.5) {Search\\检索Agent};
        \node[agent] (screen)  at (2.8,1.5) {Screen\\筛选Agent};
        \node[agent] (cluster) at (5.6,1.5) {Cluster\\聚类Agent};
        \node[agent] (summary) at (8.4,1.5) {Summary\\摘要Agent};

        % 工作区
        \node[rectangle, draw, rounded corners, fill=green!10, minimum width=10cm, minimum height=0.8cm, align=center, font=\small] (ws) at (4.2,0) {Shared Workspace：文献 $\cdot$ 嵌入 $\cdot$ 聚类 $\cdot$ 报告};

        % 箭头：用户→协调者
        \draw[arr] (input) -- (coord);

        % 箭头：协调者→第一个Agent（链式，不扇形展开避免穿模）
        \draw[arr] (coord.south) -- ++(0,-0.4) -| (search.north);

        % 箭头：Agent顺序传递
        \draw[arr] (search) -- (screen);
        \draw[arr] (screen) -- (cluster);
        \draw[arr] (cluster) -- (summary);

        % 箭头：协调者→RQ管理器
        \draw[darr] (coord) -- (rq);

        % 箭头：各Agent→工作区
        \draw[arr] (search.south)  -- (ws.north -| search);
        \draw[arr] (screen.south)  -- (ws.north -| screen);
        \draw[arr] (cluster.south) -- (ws.north -| cluster);
        \draw[arr] (summary.south) -- (ws.north -| summary);

        % 箭头：RQ管理器→工作区（绕过summary，走右侧外边）
        \draw[darr] (rq.east) -- ++(0.4,0) |- ([xshift=0.5cm]ws.north east) -- (ws.north east);
    \end{tikzpicture}
    \caption{多Agent文献综述系统架构。Coordinator链式编排四个Agent，RQ Manager提供层级研究问题驱动筛选、聚类和摘要阶段。}
    \label{fig:architecture}
\end{figure}

\subsection{流水线阶段}

处理流水线由四个顺序阶段组成，每个阶段由专门的Agent负责：

\begin{table}[htbp]
    \centering
    \caption{流水线阶段及其功能}
    \label{tbl:pipeline}
    \begin{tabular}{llll}
    \toprule
    阶段 & Agent & 核心功能 & 关键技术 \\
    \midrule
    检索 & SearchAgent  & 多数据库文献检索 & arXiv, PubMed, DBLP, OpenAlex \\
    筛选 & ScreenAgent  & 相关性过滤 & BGE嵌入, 余弦相似度 \\
    聚类 & ClusterAgent & 主题发现 & GPU-PCA, HDBSCAN \\
    摘要 & SummaryAgent & 报告生成 & 两阶段LLM调用 \\
    \bottomrule
    \end{tabular}
\end{table}

\subsection{层级研究问题}

本系统的一项关键设计是使用层级研究问题（RQ）树驱动整个流水线。RQ树包含三个层级：

\begin{itemize}
    \item \textbf{第一级（宏观维度）}：方法论、应用、趋势等大方向，定义综述的整体范围。
    \item \textbf{第二级（中观分析）}：每个维度下的具体分析角度，如方法类型、应用领域、演进时间线。
    \item \textbf{第三级（微观深入）}：技术细节，包括核心数学公式与推导链条、第一性原理理论极限、模型假设与有效性边界。
\end{itemize}

RQ树根据研究主题自动生成，以JSON文件持久化。它为检索阶段提供上下文查询、为筛选阶段提供相关性标准、为聚类阶段提供语义锚点、为摘要阶段提供结构化Prompt。

\subsection{BaseAgent抽象}

所有Agent继承自通用的\texttt{BaseAgent}抽象基类，该基类提供：
\begin{itemize}
    \item LLM调用管理，通过全局\texttt{asyncio.Semaphore(5)}进行速率限制
    \item 标准化结果格式（\texttt{AgentResult}），包含成功/错误/指标字段
    \item 工作区集成，支持数据存储和检索
    \item 结构化日志，支持进度追踪
    \item 模板方法模式：子类实现\texttt{execute()}方法，基类负责编排
\end{itemize}

\subsection{共享工作区}

\texttt{SharedWorkspace}提供两级隔离的持久化存储：
\begin{itemize}
    \item \textbf{全局共享}：文献记录和嵌入向量在所有研究主题间共享，支持复用。
    \item \textbf{主题隔离}：聚类、摘要和报告按研究主题隔离。
    \item \textbf{检查点支持}：工作区支持创建和恢复检查点，实现工作流中断和恢复。
\end{itemize}
